\section{Introdução}

O aumento contínuo do número de vulnerabilidades divulgadas em sistemas computacionais representa um dos principais desafios para a área de Segurança da Informação. Organizações de diferentes setores precisam lidar diariamente com um grande volume de vulnerabilidades, tornando inviável a correção imediata de todas elas. Nesse contexto, a priorização de remediação torna-se uma atividade fundamental para reduzir riscos e direcionar recursos para as vulnerabilidades com maior potencial de causar impactos reais.

Embora métricas tradicionais, como o Common Vulnerability Scoring System (CVSS), sejam amplamente utilizadas para avaliar a severidade das vulnerabilidades, elas nem sempre refletem a probabilidade de exploração em ambientes reais. Diversos fatores adicionais, como a disponibilidade de exploits públicos, a exploração ativa por grupos maliciosos, a presença de vulnerabilidades em campanhas recentes de ataques e informações provenientes de fontes de inteligência de ameaças, podem influenciar significativamente o risco associado a uma vulnerabilidade.

A utilização de técnicas de análise de dados e aprendizado de máquina tem se mostrado uma alternativa promissora para apoiar a identificação de vulnerabilidades com maior probabilidade de exploração. Entretanto, o desenvolvimento desses modelos depende da disponibilidade de conjuntos de dados abrangentes, atualizados e integrados, capazes de reunir informações provenientes de diferentes fontes de inteligência de vulnerabilidades. Apesar da existência de diversas bases públicas, os dados frequentemente encontram-se dispersos, heterogêneos e com diferentes formatos de representação, dificultando sua utilização em pesquisas e aplicações práticas.

Diante desse cenário, este trabalho tem como objetivo construir uma base integrada de inteligência de vulnerabilidades destinada a apoiar estudos de predição de exploração ativa e priorização de remediação. Para isso, serão coletados, processados e integrados dados provenientes de múltiplas fontes públicas de vulnerabilidades e ameaças, buscando consolidar informações relevantes em um único dataset estruturado.